\begin{tikzpicture}[font=\small]
  % ---------- chaining ----------
  \begin{scope}
    \node[anchor=south, font=\small] at (3.4,3.6){链地址法：哈希把\textbf{大宇宙}压进\textbf{小表}，冲突用链表兜底};

    % universe
    \draw[alggray, line width=0.9pt, fill=alglight!10, rounded corners=4pt] (-2.6,0.0) rectangle (-0.4,3.0);
    \node[note, anchor=south] at (-1.5,3.05){键的宇宙 $U$};
    \foreach \y/\k in {2.5/$k_1$, 1.9/$k_2$, 1.3/$k_3$, 0.7/$k_4$}
      \node[font=\footnotesize] at (-1.5,\y) {\k};

    % table slots
    \foreach \i in {0,...,4}
      \node[cell, minimum width=1.0cm] (slot\i) at (1.2,2.6-\i*0.7) {};
    \foreach \i/\lb in {0/$0$, 1/$1$, 2/$2$, 3/$3$, 4/$4$}
      \node[note, anchor=east] at (0.6,2.6-\i*0.7) {\lb};
    \node[note, anchor=south] at (1.2,3.05){表 $T$，大小 $m$};

    % hash arrows
    \draw[dedge, algblue] (-1.15,2.5) -- (0.68,2.6);
    \draw[dedge, algblue] (-1.15,1.9) -- (0.68,1.2);
    \draw[dedge, algblue] (-1.15,1.3) -- (0.68,1.2);
    \draw[dedge, algblue] (-1.15,0.7) -- (0.68,0.5);
    \node[algblue, font=\footnotesize, anchor=south] at (-0.1,2.75){$h(k)$};

    % chains
    \node[cell, minimum width=0.85cm, fill=algblue!10] (c1) at (2.5,2.6) {$k_1$};
    \draw[dedge] (slot0) -- (c1);

    \node[cell, minimum width=0.85cm, fill=algred!10] (c2) at (2.5,1.2) {$k_2$};
    \node[cell, minimum width=0.85cm, fill=algred!10] (c3) at (3.6,1.2) {$k_3$};
    \draw[dedge] (slot2) -- (c2);
    \draw[dedge] (c2) -- (c3);
    \node[algred, font=\footnotesize, anchor=west] at (4.3,1.2){冲突 $\Rightarrow$ 串成链};

    \node[cell, minimum width=0.85cm, fill=algblue!10] (c4) at (2.5,0.5) {$k_4$};
    \draw[dedge] (slot3) -- (c4);

    \node[note, anchor=north, align=center] at (2.4,-0.5)
      {$\lvert U\rvert$ 可以是 $2^{64}$，$m$ 只有几千 —— 冲突\textbf{必然}发生（抽屉原理），问题只是「平均多长」。};
  \end{scope}

  % ---------- why expected O(1) ----------
  \begin{scope}[yshift=-3.4cm]
    \node[anchor=south, font=\small] at (5.0,2.2){为什么期望是 $O(1)$};

    \node[blk, minimum width=3.4cm, fill=algteal!7, draw=algteal, anchor=north west] (h1) at (0,1.4)
      {\textbf{全域哈希}\\[2pt]\footnotesize 从一族 $\mathcal{H}$ 中\textbf{随机}选 $h$\\
       \footnotesize 使 $\Pr[h(k_i)=h(k_j)]\le 1/m$};
    \node[blk, minimum width=3.4cm, anchor=north west] (h2) at (4.0,1.4)
      {\textbf{链长期望}\\[2pt]\footnotesize $E[\text{链长}]=1+\alpha$\\
       \footnotesize 装载因子 $\alpha=n/m$};
    \node[blk, minimum width=3.4cm, fill=algblue!7, draw=algblue, anchor=north west] (h3) at (8.0,1.4)
      {\textbf{保持 $\alpha=O(1)$}\\[2pt]\footnotesize 表满了就\textbf{翻倍重哈希}\\
       \footnotesize 摊还后仍是 $O(1)$};
    \draw[dedge] (3.5,0.5) -- (4.0,0.5);
    \draw[dedge] (7.5,0.5) -- (8.0,0.5);

    \node[note, anchor=north west, align=left] at (0,-0.6)
      {\textbf{随机性来自算法而不是数据} —— 这一点很关键。若固定一个 $h$，对手总能构造出全部撞在\\
       同一个槽的输入，把查找退化成 $O(n)$。随机选 $h$ 之后，「最坏输入」这个概念就失效了。\\[4pt]
       代价：$O(1)$ 只是\textbf{期望}值，不是最坏保证。要最坏 $O(\log n)$ 保证就得用平衡树 —— 而树还额外\\
       支持 \texttt{find\_prev} / \texttt{find\_next} / 按序遍历，哈希表\textbf{做不到}。这是选型时真正的分界线。};
  \end{scope}
\end{tikzpicture}
